\documentclass[11pt]{article}

\usepackage[margin=2.5cm]{geometry}
\usepackage{booktabs}
\usepackage{longtable}
\usepackage{array}
\usepackage{listings}
\usepackage{xcolor}
\usepackage{hyperref}
\usepackage{parskip}

\hypersetup{
    colorlinks=true,
    linkcolor=blue,
    urlcolor=blue,
    citecolor=blue
}

\lstdefinestyle{Rstyle}{
    language=R,
    basicstyle=\ttfamily\small,
    keywordstyle=\color{blue},
    commentstyle=\color{gray},
    stringstyle=\color{teal},
    numbers=left,
    numberstyle=\tiny\color{gray},
    breaklines=true,
    frame=single,
    backgroundcolor=\color{gray!5},
    showstringspaces=false
}

\title{A Cleaned Global Earthquake Dataset for AddiVortes Benchmarking}
\author{}
\date{\today}

\begin{document}

\maketitle

\section{Introduction}

This article documents a cleaned dataset of global earthquake events prepared for
benchmarking the AddiVortes model. The dataset was chosen because it provides a
large number of observations with a genuine mix of numerical covariates measured
on fundamentally different scales, together with categorical covariates, and
latitude/longitude coordinates spanning the entire globe. The latter is
particularly useful for testing distance-based splitting rules, since great-circle
distance is the natural metric for the coordinate pair while Euclidean distance
is appropriate for the remaining numerical covariates.

\section{Data Source}

The raw data was obtained from the U.S. Geological Survey's Earthquake Hazards
Program, via the public ANSS Comprehensive Earthquake Catalog (ComCat) query API:

\begin{center}
\url{https://earthquake.usgs.gov/earthquakes/search/}
\end{center}

The pull used for this article consists of the most recent 20{,}000 earthquake
records in the catalog which occured in 2025 (magnitude $\geq 2.5$), which
corresponds to events occurring between 5 May 2025 and 31 December 2025. The USGS
API caps any single query at 20{,}000 returned rows; a full multi-decade pull
requires paginating over the \texttt{offset} parameter or splitting the request
into shorter date ranges, but the analysis in this article is restricted to this
single 20{,}000-row pull as a representative working sample.

The raw file, \texttt{raw\_data.csv}, contained 20{,}000 rows and 19 columns:
\texttt{time}, \texttt{latitude}, \texttt{longitude}, \texttt{depth}, \texttt{mag},
\texttt{magType}, \texttt{nst}, \texttt{gap}, \texttt{dmin}, \texttt{rms},
\texttt{net}, \texttt{type}, \texttt{horizontalError}, \texttt{depthError},
\texttt{magError}, \texttt{magNst}, \texttt{status}, \texttt{locationSource}, and
\texttt{magSource}.

\section{Data Cleaning}

The raw data required six cleaning steps: (i) restricting to genuine earthquake
events, since the catalog also contains a small number of mining explosions and
other non-tectonic events; (ii) removing columns that were pure metadata,
near-duplicates of another retained column, or deliberately excluded from the
benchmark covariate set; (iii) removing the small proportion of rows with
residual missing values in the quality-control columns; (iv) subsetting to
reviewed body-wave magnitude events, i.e. \texttt{magType == "mb"} and
\texttt{status == "reviewed"}; and (v) dropping the now-constant
\texttt{magType} and \texttt{status} columns from the retained covariate set.
The full cleaning pipeline, implemented in R, is reproduced below.

\begin{lstlisting}[style=Rstyle, caption={R script used to clean the raw USGS earthquake catalog extract.}, label={lst:clean}]
# =========================================================
# AddiVortes test data prep: USGS earthquake catalog cleanup
# =========================================================

# ---- 1. Read in the raw CSV ----
df <- read.csv("raw_data.csv", stringsAsFactors = FALSE)

cat("Starting rows:", nrow(df), "\n")
cat("Starting columns:", ncol(df), "\n\n")

# ---- 2. Filter to real earthquakes only ----
# 'type' includes a small number of mining explosions, landslides, etc.
# which are noise for an earthquake-magnitude model.
before_n <- nrow(df)
df <- df[df$type == "earthquake", ]
cat("After filtering to type == 'earthquake':", nrow(df),
    "rows (dropped", before_n - nrow(df), ")\n")

# ---- 3. Drop redundant / non-covariate columns ----
# - 'type' is now constant (all 'earthquake') -> no information left
# - 'locationSource' is ~99.99% identical to 'net' -> redundant
# - 'magSource' is ~99.7% identical to 'net' -> redundant
# - 'time' dropped for this non-time-series test (re-add + engineer
#    year/month/hour if you want temporal covariates later)
# - 'gap' removed from the benchmark covariate set as requested
drop_cols <- c("type", "locationSource", "magSource", "time", "gap")
df <- df[, !(names(df) %in% drop_cols)]
cat("Dropped columns:", paste(drop_cols, collapse = ", "), "\n")
cat("Remaining columns:", paste(names(df), collapse = ", "), "\n\n")

# ---- 4. Drop rows with any remaining missing values ----
# Missingness across the numeric quality/error columns is low (~5.6%)
# and concentrated in a subset of 'ml' and 'mw' events, so row deletion
# is cheap and keeps the dataset fully complete.
before_n <- nrow(df)
df_clean <- df[complete.cases(df), ]
dropped_n <- before_n - nrow(df_clean)
cat("After dropping rows with missing values:", nrow(df_clean),
    "rows (dropped", dropped_n,
    sprintf("[%.1f%%]", 100 * dropped_n / before_n), ")\n\n")

# ---- 5. Restrict to the benchmark subset used in this session ----
# Keep only reviewed body-wave magnitude measurements for the benchmark.
before_n <- nrow(df_clean)
df_clean <- df_clean[df_clean$magType == "mb" & df_clean$status == "reviewed", ]
cat("After filtering to magType == 'mb' and status == 'reviewed':",
    nrow(df_clean), "rows (dropped", before_n - nrow(df_clean), ")\n\n")

# ---- 6. Drop now-constant filtering columns ----
# After subsetting, 'magType' and 'status' are constant and carry no information.
df_clean <- df_clean[, !(names(df_clean) %in% c("magType", "status"))]
cat("After dropping constant columns magType and status:",
    nrow(df_clean), "rows and", ncol(df_clean), "columns\n\n")

# ---- 7. Sanity checks ----
cat("Final shape:", nrow(df_clean), "rows x", ncol(df_clean), "columns\n")
cat("Total missing values remaining:", sum(is.na(df_clean)), "\n\n")

cat("Categorical cardinality:\n")
for (col in c("net")) {
  cat(" ", col, "->", length(unique(df_clean[[col]])), "levels\n")
}

# ---- 8. Write out the cleaned file ----
write.csv(df_clean, "earthquakes.csv", row.names = FALSE)
cat("\nSaved cleaned file to earthquakes.csv\n")
\end{lstlisting}

Applying this pipeline reduced the raw extract from 20{,}000 rows and 19 columns
to a final cleaned dataset of 11{,}697 rows and 12 columns, with zero remaining
missing values. Table~\ref{tab:pipeline} summarises each stage.

\begin{table}[h]
\centering
\begin{tabular}{lrr}
\toprule
Stage & Rows & Columns \\
\midrule
Raw extract & 20{,}000 & 19 \\
After filtering to \texttt{type == "earthquake"} & 19{,}841 & 19 \\
After dropping redundant columns & 19{,}841 & 14 \\
After dropping incomplete rows & 18{,}732 & 14 \\
After subsetting to \texttt{magType == "mb"} and \texttt{status == "reviewed"} & 11{,}697 & 14 \\
After dropping constant \texttt{magType} and \texttt{status} columns & 11{,}697 & 12 \\
\bottomrule
\end{tabular}
\caption{Row and column counts at each stage of the cleaning pipeline.}
\label{tab:pipeline}
\end{table}

\section{Covariate Descriptions}

Table 1 describes each of the 12 covariates retained in the
final cleaned dataset, along with its type and the range (for numeric variables)
or levels (for categorical variables) observed in the data.

\begin{longtable}{p{2.6cm}p{1.6cm}p{6.3cm}p{3.5cm}}
\toprule
\textbf{Covariate} & \textbf{Type} & \textbf{Description} & \textbf{Range / Levels} \\
\midrule
\endfirsthead
\toprule
\textbf{Covariate} & \textbf{Type} & \textbf{Description} & \textbf{Range / Levels} \\
\midrule
\endhead

\texttt{latitude} & Numeric & Epicenter latitude in decimal degrees (WGS84). Natural great-circle covariate. & $-65.6^{\circ}$ to $87.1^{\circ}$ \\
\addlinespace
\texttt{longitude} & Numeric & Epicenter longitude in decimal degrees (WGS84). Natural great-circle covariate. & $-180.0^{\circ}$ to $180.0^{\circ}$ \\
\addlinespace
\texttt{depth} & Numeric & Depth of the hypocenter below the surface, in kilometers. Small negative values reflect very shallow events located slightly above the reference datum. & $-3.4$ to $669.6$ km \\
\addlinespace
\texttt{mag} & Numeric & Earthquake magnitude, on the body-wave scale used for the retained events. Natural target variable. & $2.5$ to $8.8$ \\
\addlinespace
\texttt{nst} & Numeric & Number of seismic stations used to determine the event's \emph{location}. & $3$ to $566$ \\
\addlinespace
\texttt{dmin} & Numeric & Horizontal distance, in degrees, from the epicenter to the nearest reporting station. & $0$ to $55.9$ \\
\addlinespace
\texttt{rms} & Numeric & Root-mean-square of the travel-time residuals, in seconds; a location goodness-of-fit measure. & $0$ to $2.0$ s \\
\addlinespace
\texttt{net} & Categorical & Seismic network that contributed the preferred origin (e.g.\ national or regional networks). & 12 levels, dominated by \texttt{us} (16{,}088 rows) \\
\addlinespace
\texttt{horizontalError} & Numeric & 95\% confidence radius on the epicenter location, in kilometers. & $0$ to $40.8$ km \\
\addlinespace
\texttt{depthError} & Numeric & 95\% confidence uncertainty on the depth estimate, in kilometers. & $0$ to $42.9$ km \\
\addlinespace
\texttt{magError} & Numeric & Standard error on the magnitude estimate, in magnitude units. & $0$ to $1.3$ \\
\addlinespace
\texttt{magNst} & Numeric & Number of seismic stations used to compute the \emph{magnitude} (distinct from \texttt{nst}). & $0$ to $1{,}027$ \\
\addlinespace
\bottomrule
\caption{Covariates in the cleaned earthquake dataset.}
\label{tab:covariates}
\end{longtable}

\section{Notes for Model Testing}

One covariate warrants caution during testing: \texttt{net} is highly
imbalanced, with the \texttt{us} network accounting for roughly 80\% of
observations; this may limit the effective diversity of that categorical
covariate for split-finding. Aside from that variable, the dataset offers a
genuine mix of scales and types --- degrees, kilometers, seconds, magnitude
units, and station counts among the numeric covariates, together with a single
categorical covariate --- with global latitude/longitude coverage and no missing
data, making it well suited to benchmarking AddiVortes against both great-circle
and Euclidean distance metrics.

\end{document}
